\chapter{Building a New Simulation}
\label{chap:new-simulation}

This chapter specifies how a scenario is assembled in \lupnt{}: the object model
(\cpp{Simulation}, \cpp{World}, \cpp{Agent}, \cpp{Application}, \cpp{Device},
\cpp{Measurement}, \cpp{Dynamics}, \cpp{State}), the discrete-event engine that
schedules them, the YAML schema that instantiates them through a string-keyed
asset factory, and the two supported authoring paths --- a factory-registered
C++ class, or a pure-Python subclass injected through the \code{pybind11}
trampolines. The implementing sources are
\file{cpp/lupnt/simulations/simulation.h}, \file{cpp/lupnt/simulations/world.h},
\file{cpp/lupnt/core/event.h}, \file{cpp/lupnt/core/asset_factory.h},
\file{cpp/lupnt/agents/agent.h}, \file{cpp/lupnt/applications/application.h},
\file{cpp/lupnt/measurements/measurement.h}, the shipped scenarios under
\file{configs/}, and the Python binding layer
\file{python/bindings/py_simulation.cc}. The design follows the agent-based
architecture described in \citep{iiyama2023lupnt,casadesus2025lupnt} and
\citep{iiyama2026thesis}.

\section{One engine, one path}
\label{sec:sim-one-engine}

Every scenario in \lupnt{} is built from a YAML configuration by the same generic
event-driven engine, \cpp{lupnt::Simulation}
(\file{cpp/lupnt/simulations/simulation.cc}). There is no ``monolithic''
alternative: nothing in the repository subclasses \cpp{Simulation} or hand-writes
a \cpp{Run()} loop. Agents, applications, devices, dynamics, channels, and
filters are all assembled by name through factories; a top-level \code{world:}
block defines the shared read-only environment (epoch, integration frame, force
model) that both the truth agents and the estimators draw from; and any new
mission or navigation logic goes into a factory-registered \cpp{Application} that
self-drives from that environment.

The ground-station ODTS, GNSS ODTS, inter-satellite-link (ISL) ODTS (centralized
\emph{and} distributed), ephemeris, surface-rover, and lander scenarios all run
this one way; see \file{configs/ground_station_odts.yaml},
\file{configs/lunar_gnss_odts.yaml}, \file{configs/isl_odts_distributed.yaml},
\file{configs/ephemeris.yaml}, \file{configs/surface_rover_nav.yaml},
\file{configs/lander_nav.yaml}, and \file{configs/sat_bearing_odts.yaml}.

A configuration is launched from Python with \py{pnt.Simulation(yaml\_or\_dict)}
--- the pattern used by every \file{python/examples/exN_run_*.py} driver script
--- or from a small C++ driver such as the tutorials under
\file{cpp/examples/tutorials/}, which merely load the YAML, construct a
\cpp{Simulation}, and call \cpp{Run()}:

\begin{lstlisting}[language=cpp, caption={The complete C++ driver for a scenario.}, label={lst:sim-cpp-driver}]
// cpp/examples/tutorials/ex7_groundstation_odts.cc :: main
YAML::Node config = YAML::LoadFile(config_path);
Simulation sim(config);          // Setup() runs inside the constructor
sim.Run();                       // drains the event queue

// The manager agent's application owns the centralized OD solution.
Agent* mgr_agent = sim.GetAgent("gs_manager");
auto mgr = std::dynamic_pointer_cast<GroundStationManagerApp>(mgr_agent->GetApplication());
LUPNT_CHECK(mgr && mgr->HasSolved(), "GroundStationManagerApp did not solve", "ex7");
\end{lstlisting}

The only real design choice left to the scenario author is \emph{how to decompose
the per-epoch logic across applications}. Four shapes recur, and every shipped
scenario is one of them:

\begin{description}[leftmargin=1.2cm,style=nextline,font=\normalfont\itshape]
  \item[Sensor / estimator split] Separate sensor applications generate
    observations and push them to a manager application that estimates. Ground
    station OD: \cpp{GroundStationTrackingApp} on each \cpp{GroundStation}, all
    feeding one \cpp{GroundStationManagerApp} on a \cpp{GroundStationManager}
    agent.
  \item[Single coordinator application] One manager agent hosts an application
    that runs the whole estimator itself. Centralized ground ODTS:
    \cpp{SurfaceStationManager} $+$ \cpp{GroundOdtsApp}, a single EKF over all
    surface-station beacon observations.
  \item[Fully distributed] Every platform is its own agent running its own
    application, exchanging state over the \cpp{Publish} / \cpp{Subscribe} bus.
    Distributed ISL ODTS: \cpp{Spacecraft} $+$ \cpp{SatelliteOdtsApp}, alongside
    \cpp{SurfaceStation} $+$ \cpp{StationBeaconSensor} feeding
    \cpp{SurfaceStationManager} $+$ \cpp{GroundOdtsApp}.
  \item[Several applications on one platform] One physical agent hosts multiple
    applications in an \code{applications:} list, run in insertion order and
    sharing the agent's truth state. Lander descent: a guidance
    \cpp{LanderGncApp} that owns the truth trajectory, plus a navigation
    \cpp{LanderNavApp} multiplicative extended Kalman filter (MEKF) that reads it,
    resolved through \cpp{GetApplicationByName}.
\end{description}

A heavyweight engine needs no rewrite to fit. From its \cpp{Step}, the
application \cpp{LunarGnssOdtsApp} calls
\cpp{RunLunarGnssODTSMonteCarlo} directly, so an existing batch or
Monte-Carlo solver can be dropped inside a single application unchanged.

\section{The object model}
\label{sec:sim-object-model}

\Cref{tab:sim-blocks} lists the eight building blocks and the header that defines
each one; \cref{fig:sim-object-model} shows containment and the per-step call
flow.

\begin{table}[H]
  \centering
  \caption{The building blocks of a \lupnt{} scenario. Paths are relative to
           \texttt{cpp/lupnt/}.}
  \label{tab:sim-blocks}
  \begin{tabularx}{\textwidth}{@{}l L L@{}}
    \toprule
    Block & Base class (header) & Role \\
    \midrule
    Simulation & \cpp{Simulation} (\file{simulations/simulation.h}) &
      Discrete-event engine: a time-ordered \cpp{std::priority\_queue<Event>},
      \cpp{Schedule} / \cpp{Run} / publish--subscribe, and the \cpp{agents\_},
      \cpp{channels\_}, \cpp{constellations\_} registries built from config. Owns
      the \cpp{World}. \\
    World & \cpp{World} (\file{simulations/world.h}) &
      The shared, read-only physical environment built from the top-level
      \code{world:} block: epoch, integration \code{frame}, one force model,
      an optional point-mass \code{gravity:}, an optional \code{dem:} terrain
      service, and an optional \code{plasma:} block. \cpp{MakeDynamics()} hands
      out a fresh \cpp{NBodyDynamics} from that one force model. It never
      propagates agents. \\
    Agent & \cpp{Agent} $\rightarrow$ \cpp{AgentWithDynamics}
      (\file{agents/agent.h}) &
      A platform (satellite, ground station, rover, lander, surface station).
      Owns a \cpp{Dynamics}, \cpp{Device}s, a \cpp{State}, and hosts one or more
      \cpp{Application}s. Answers \cpp{GetStateAt(t)} for measurement geometry. \\
    Application & \cpp{Application} (\file{applications/application.h}) &
      The mission / navigation logic run per step. Builds \cpp{Measurement}s and
      runs a \cpp{Filter}. \\
    Device & \cpp{Device} (\file{devices/device.h}) &
      A sensor or component on an agent (\cpp{Clock}, \cpp{Imu}, \cpp{Camera},
      \cpp{Transmitter} / \cpp{Receiver} / \cpp{Transponder},
      \cpp{GnssTransmitter} / \cpp{GnssReceiver}). Steps on its own schedule; read
      by applications through \cpp{GetDevice(name)}. \\
    Measurement & \cpp{Measurement} / \cpp{ErrorStateMeasurement}
      (\file{measurements/measurement.h}) &
      The observable model. \cpp{Compute(x, H)} returns the predicted value, its
      covariance, and the Jacobian $\mat{H} = \partial \vec{h} / \partial \vec{x}$;
      \cpp{CreateFunction()} wires it into a \cpp{Filter}. \\
    Dynamics & \cpp{Dynamics} (\file{dynamics/dynamics.h}) &
      The propagation model. \cpp{Propagate(x0,t0,tf[,u][,stm])} advances a
      \cpp{State} and optionally its state-transition matrix
      $\mat{\Phi} = \partial \vec{x}_f / \partial \vec{x}_0$. \\
    State & \cpp{State}, \cpp{JointState} (\file{states/state.h},
      \file{states/joint_state.h}) &
      A labeled Eigen vector with a type, per-element names and units, and a
      frame. \cpp{JointState} composes orbit $+$ clock ($+$ parameters) into one
      filter state. \\
    \bottomrule
  \end{tabularx}
\end{table}

\begin{figure}[H]
  \centering
  \begin{tikzpicture}[node distance=9mm and 12mm]
    % --- engine layer -------------------------------------------------------
    \node[boxnode, minimum width=3.6cm] (sim)
      {\textbf{Simulation}\\[1pt]\scriptsize priority queue of Events\\
       \scriptsize agents, channels, constellations};
    \node[boxnode, left=16mm of sim, minimum width=3.0cm] (world)
      {\textbf{World}\\[1pt]\scriptsize epoch, frame, force model\\
       \scriptsize gravity, DEM, plasma};

    % --- agent layer --------------------------------------------------------
    \node[boxnode, below=16mm of sim, xshift=-26mm, minimum width=3.0cm] (agentA)
      {\textbf{Agent} (truth)\\[1pt]\scriptsize Dynamics, Devices, State};
    \node[boxnode, below=16mm of sim, xshift=26mm, minimum width=3.0cm] (agentB)
      {\textbf{Agent} (host)\\[1pt]\scriptsize Devices, State};

    % --- application layer --------------------------------------------------
    \node[boxnode, below=14mm of agentB, minimum width=3.4cm] (app)
      {\textbf{Application}\\[1pt]\scriptsize Setup / Step / Log};

    % --- model layer --------------------------------------------------------
    \node[boxnode, below=14mm of app, xshift=-24mm, minimum width=2.6cm] (meas)
      {\textbf{Measurement}\\[1pt]\scriptsize $z=h(x)$, $H$, $R$};
    \node[boxnode, below=14mm of app, xshift=24mm, minimum width=2.6cm] (filt)
      {\textbf{Filter}\\[1pt]\scriptsize EKF / UKF / SRIF / batch};
    \node[boxnode, below=12mm of meas, xshift=24mm, minimum width=3.4cm] (state)
      {\textbf{JointState} $+$ \textbf{Dynamics}\\[1pt]\scriptsize $\Phi$, process noise};

    % --- edges --------------------------------------------------------------
    \draw[flow] (sim) -- node[above,font=\scriptsize\sffamily] {owns} (world);
    \draw[flow] (sim) -- node[left,pos=0.55,font=\scriptsize\sffamily] {Step (AGENT)} (agentA);
    \draw[flow] (sim) -- node[right,pos=0.55,font=\scriptsize\sffamily] {Step (AGENT)} (agentB);
    \draw[flow] (agentB) -- node[right,font=\scriptsize\sffamily] {hosts} (app);
    \draw[flow] (app) -- node[left,font=\scriptsize\sffamily] {builds} (meas);
    \draw[flow] (app) -- node[right,font=\scriptsize\sffamily] {drives} (filt);
    \draw[flow] (meas) -- (state);
    \draw[flow] (filt) -- (state);
    \draw[flow, dashed] (app.west) to[out=170,in=-100]
      node[left,pos=0.6,font=\scriptsize\sffamily] {GetStateAt$(t)$} (agentA.south);
    \draw[flow, dashed] (world.south) to[out=-90,in=170]
      node[left,pos=0.35,font=\scriptsize\sffamily] {MakeDynamics} (agentA.west);
  \end{tikzpicture}
  \caption{Object model and per-step call flow. Solid arrows are ownership or
           scheduling; dashed arrows are read-only queries. The engine pops due
           events, agents self-propagate their truth, devices advance, and the
           hosted application reads truth through \cpp{GetStateAt}, forms a
           measurement, and runs one filter cycle.}
  \label{fig:sim-object-model}
\end{figure}

\subsection{Simulation --- the engine}
\label{sec:sim-engine}

\cpp{Simulation} is a discrete-event engine backed by a
\cpp{std::priority\_queue<Event>}. An \cpp{Event} (\file{cpp/lupnt/core/event.h})
is a \code{\{time, frequency, priority, callback\}} tuple:

\begin{lstlisting}[language=cpp, caption={The event record and its ordering.}, label={lst:sim-event}]
// cpp/lupnt/core/event.h :: Event
struct Event {
  Real time;
  Real frequency;
  Real priority = 0.0;
  std::function<void(Real)> callback;
  bool operator<(const Event& e) const {
    return time > e.time || (time == e.time && priority > e.priority);
  }
  enum Priority {
    LOW = 0, MEDIUM = 1, HIGH = 2,
    LOGGING = LOW, APPLICATION = LOW,
    DEVICE = MEDIUM,
    DYNAMICS = HIGH, AGENT = HIGH,
  };
  constexpr static Real SINGLE_EVENT = -1.0;
};
\end{lstlisting}

The comparator inverts \cpp{operator<} on \code{time}, so the
\cpp{std::priority\_queue} (a max-heap) pops the \emph{earliest} event first ---
the intended behaviour. The tie-break on \code{priority} is written in the same
direction, so at equal time the numerically \emph{lower} priority value is popped
first: \code{APPLICATION} $=$ \code{LOGGING} $(0)$, then \code{DEVICE} $(1)$, then
\code{DYNAMICS} $=$ \code{AGENT} $(2)$.

\begin{lupntnote}
The enum names suggest the opposite intent (dynamics before devices before
applications). In practice the ordering is benign, because an application never
reads a cached agent state: \cpp{AgentWithDynamics::GetStateAt(t)} propagates a
\emph{copy} of \code{state\_} from \code{time\_} to $t$ without mutating the
agent, so the truth an application sees at time $t$ is correct whether or not the
agent's own \cpp{Step} has already fired at that epoch. Applications that depend
on a device having advanced first (an IMU sample, a clock realization) should
read the device through its own accessor rather than relying on event order.
\end{lupntnote}

The main loop pops the earliest event, sets the simulation clock, invokes the
callback, and --- if the event carries a non-zero \code{frequency} --- reschedules
it at $t + 1/f$, which is how periodic \cpp{Step} callbacks recur:

\begin{lstlisting}[language=cpp, caption={The engine loop.}, label={lst:sim-run-loop}]
// cpp/lupnt/simulations/simulation.cc :: Simulation::Run
running_ = true;
while (!queue_.empty() && queue_.top().time <= duration_ + EPS) {
  Event e = queue_.top();
  queue_.pop();
  time_ = e.time;
  e.callback(e.time);                 // Agent / Device / Application Step, or a Publish
  if (e.frequency > 0.0) {            // periodic -> reschedule
    Real time_next = e.time + 1.0 / e.frequency;
    if (time_next <= duration_ + EPS) { e.time = time_next; queue_.push(e); }
  }
}
DataLogger::Flush();
\end{lstlisting}

\cpp{Schedule(time, callback, freq, priority)} adds events. \cpp{Subscribe(topic,
callback)} and \cpp{Publish(time, topic, message)} provide a topic-based message
bus; a \cpp{Publish} does not call subscribers inline but schedules one
single-shot \code{APPLICATION}-priority event per subscriber at the given time,
which is what gives the distributed ISL scenario its one-interval message latency.
\cpp{GetAgent(name)} and \cpp{GetChannel(name)} accept either a bare name or the
fully-qualified \code{<simname>/<name>} key under which the object was registered.

\subsection{World --- the shared environment}
\label{sec:sim-world}

\cpp{Simulation::Setup} builds a single \cpp{World} from the top-level
\code{world:} block \emph{after} the epoch is set and \emph{before} any agents, so
every agent and application can pull the shared force model and truth facade
during its own \cpp{Setup}. Reach it from anywhere with \cpp{Agent::GetWorld()}
(equivalently \cpp{GetSimulation()->GetWorld()}).

The \cpp{World} recognizes the keys in \cref{tab:sim-world-keys}. All of them are
optional; a scenario with no \code{world:} block simply has
\cpp{GetWorld() == nullptr}.

\begin{table}[H]
  \centering
  \caption{Keys parsed by \cpp{World::World} in \texttt{cpp/lupnt/simulations/world.cc}.}
  \label{tab:sim-world-keys}
  \begin{tabularx}{\textwidth}{@{}l l L@{}}
    \toprule
    Key & Default & Meaning \\
    \midrule
    \code{frame} & \code{MOON\_CI} & Integration / environment reference frame. \\
    \code{force\_model} & --- & An \cpp{NBodyDynamics}-style force-model block
      (see \cref{sec:sim-force-model}); enables \cpp{MakeDynamics()},
      \cpp{MakeTruthDynamics()} and \cpp{HasForceModel()}. \\
    \code{gravity.body} & \code{MOON} & Central body whose point-mass $GM$ backs
      \cpp{Gravity(r)}, used by surface inertial-navigation mechanization. \\
    \code{dem.site\_lat\_deg}, \code{dem.site\_lon\_deg} & required in
      \code{dem:} & Site centre of the LOLA digital elevation model, and the
      origin of the local east--north--up (ENU) tangent frame. \\
    \code{dem.half\_width\_m} & \num{5000} & Half-width of the loaded DEM tile
      [\si{\meter}]. \\
    \code{dem.max\_res\_m} & \num{20} & Target DEM resolution [\si{\meter}]. \\
    \code{dem.dem\_file} & --- & Explicit GeoTIFF to crop instead of downloading
      a tile (used by tests and small fixtures). \\
    \code{plasma} & --- & Ionosphere / plasmasphere signal-delay environment; a
      shared truth property, read by the GNSS ODTS application. \\
    \bottomrule
  \end{tabularx}
\end{table}

The epoch is not a \code{world:} key: it is taken from the global \lupnt{} epoch,
which \cpp{Simulation::Setup} sets from the scenario-level \code{epoch:} string
before constructing the \cpp{World}.

Three factory methods hand out dynamics built from the one shared force model:

\begin{lstlisting}[language=cpp, caption={\cpp{World} dynamics factories.}, label={lst:sim-world-dyn}]
// cpp/lupnt/simulations/world.cc
Ptr<NBodyDynamics> World::MakeDynamics() const {        // for ESTIMATORS
  Config cfg = force_model_;
  auto dynamics = MakePtr<NBodyDynamics>(cfg);
  dynamics->SetFrame(frame_);
  dynamics->SetAutodiff(true);      // analytic state-transition matrix
  return dynamics;
}
Ptr<NBodyDynamics> World::MakeDynamics(double cr, double area_m2, double mass_kg) const;
Ptr<NBodyDynamics> World::MakeTruthDynamics() const {   // same model, SetAutodiff(false)
  auto dynamics = MakeDynamics();
  dynamics->SetAutodiff(false);     // truth propagation needs no STM
  return dynamics;
}
\end{lstlisting}

An application that reuses \cpp{MakeDynamics()} for \emph{both} the truth grid and
the filter has no truth/filter model mismatch by construction. The \cpp{World} is
strictly read-only: it \emph{provides} the environment and a truth facade
(\cpp{World::GetStateAt(name, t)}, which forwards to
\cpp{Simulation::GetAgent(name)->GetStateAt(t)}), but each
\cpp{AgentWithDynamics} still self-propagates through its own \code{dynamics\_}.

\begin{lupntnote}
An agent that declares no \code{dynamics:} block of its own inherits the shared
\code{world.force\_model} as its truth dynamics. \cpp{Simulation::Setup} injects a
copy of that block into the agent's config, defaulting \code{class} to
\code{NBodyDynamics}, inheriting \code{world.frame}, and forcing
\code{autodiff: false}. Fixed agents (ground and surface stations, which use a
\cpp{StaticDynamics}) and manager agents ignore the injected block.
\end{lupntnote}

\subsection{Agent --- a platform}
\label{sec:sim-agent}

\cpp{Agent} owns \code{name\_}, a back-pointer to the \cpp{Simulation}, a
\cpp{State}, a map of \cpp{Device}s, and one or more \cpp{Application}s. Its
pure-virtual

\begin{equation}
  \label{eq:sim-getstateat}
  \cpp{Cart6 GetStateAt(Real t) const}
  \quad\longmapsto\quad
  \vec{x}(t) = \begin{bmatrix}\vec{r}(t)\\ \vec{v}(t)\end{bmatrix}
\end{equation}

is the contract every concrete agent must answer --- \emph{``where am I at time
$t$?''} --- and is what measurement models call for light-time geometry. The
returned state is Cartesian position and velocity in the agent's dynamics frame.

Almost every physical agent derives from \cpp{AgentWithDynamics}, which adds a
\cpp{Ptr<Dynamics>}, an \cpp{AttitudeDynamics}, a \cpp{Cart6 state\_}, an
\cpp{Attitude attitude\_}, and a \cpp{State control\_}. Its \cpp{Step(t)} calls
\cpp{Propagate(t)} --- integrating the dynamics in place from \code{time\_} to $t$
with the current \code{control\_} --- and then \cpp{Log(t)}. Its
\cpp{GetStateAt(t)} returns the cached \code{state\_} if $|t - \code{time\_}| <
\code{EPS}$, and otherwise propagates a non-mutating copy.

\Cref{tab:sim-agent-classes} lists the concrete agents registered with the
factory. Applications are attached with the \code{application:} map (a single app)
or the \code{applications:} sequence (several, run in insertion order); retrieve
them with \cpp{GetApplication()}, \cpp{GetApplications()}, or
\cpp{GetApplicationByName("...")}, which matches any app whose (agent-prefixed)
name \emph{ends with} the argument.

\begin{table}[H]
  \centering
  \caption{Agent classes registered with \cpp{AgentFactory} (the string is the
           YAML \code{class:} value).}
  \label{tab:sim-agent-classes}
  \begin{tabularx}{\textwidth}{@{}l l L@{}}
    \toprule
    \code{class:} & Source & Role \\
    \midrule
    \cpp{Satellite} & \file{agents/satellite.cc} &
      Orbit-only truth from a \code{ClassicalOE} initial state; self-propagates
      its \code{dynamics:} block. \\
    \cpp{Spacecraft} & \file{agents/spacecraft.cc} &
      Joint orbit $+$ clock 8-state truth
      $[\vec{r}, \vec{v}, b, d]$ with a \code{clock:} model; accepts
      \code{r0\_m}/\code{v0\_mps} or a \code{ClassicalOE} block. \\
    \cpp{GroundStation} & \file{agents/ground_station.cc} &
      Earth station from \code{latitude\_deg} / \code{longitude\_deg} /
      \code{altitude\_m}, or from SPICE (\code{spice\_name}, \code{spice\_id},
      \code{spice\_frame}, \code{spice\_kernel(s)}). \\
    \cpp{GroundStationManager} & \file{agents/ground_station_manager.cc} &
      Massless ground-segment host for a centralized estimator application. \\
    \cpp{SurfaceStation} & \file{agents/surface_station.cc} &
      Lunar surface beacon / station. \\
    \cpp{SurfaceStationManager} & \file{agents/surface_station_manager.cc} &
      Host for the centralized ground filter in the ISL scenario. \\
    \cpp{Rover} & \file{agents/rover.cc} &
      Surface rover driven by \cpp{SurfaceDynamics2D} with a control input. \\
    \cpp{Lander} & \file{agents/lander.cc} &
      Thin descent platform; carries the guidance and navigation applications. \\
    \cpp{EphemerisManager} & \file{agents/ephemeris_manager.cc} &
      Host for the ephemeris-generation application. \\
    \cpp{LunarNavConstellation} & \file{agents/lunar_nav_constellation.cc} &
      Group generator: one \code{agents:} entry expands into $N$ child
      \cpp{Spacecraft}, from an explicit \code{satellites:} list or a symmetric
      \code{walker:} spec. \\
    \bottomrule
  \end{tabularx}
\end{table}

\cpp{LunarNavConstellation} deserves a note because it removes the need to declare
one agent per navigation satellite. Its \code{walker:} block takes
\code{n\_planes}, \code{sats\_per\_plane}, the common frozen-orbit elements
\code{a}, \code{e}, \code{i}, \code{omega}, and the phasing keys
\code{raan0\_deg}, \code{m0\_deg}, \code{phase\_deg}, plus \code{frame} (default
\code{MOON\_OP}) and \code{name\_prefix} (default \code{SV}). Right ascension of
the ascending node for plane $p$ and mean anomaly for satellite $s$ of that plane
are

\begin{align}
  \Omega_p &= \Omega_0 + p\,\frac{360^\circ}{n_\text{planes}},
  & p &= 0,\dots,n_\text{planes}-1, \label{eq:sim-walker-raan}\\
  M_{p,s} &= M_0 + s\,\frac{360^\circ}{n_\text{sat/plane}} + p\,\Delta\phi,
  & s &= 0,\dots,n_\text{sat/plane}-1, \label{eq:sim-walker-anomaly}
\end{align}

with $\Delta\phi$ the inter-plane phasing \code{phase\_deg}. The surface-rover
example sources its relay truth from such a constellation via
\cpp{GetSatelliteStateAt(j, t)}. A separate top-level \code{constellations:} block
builds \cpp{Constellation} / \cpp{GnssConstellation} objects, which are registered
and set up alongside agents.

\subsection{Application --- the mission logic}
\label{sec:sim-application}

\cpp{Application} holds a back-pointer to its owning \cpp{Agent} and a copy of its
\code{application:} config block. The base constructor reads only two keys,
\code{name} (defaulting to the object id) and \code{frequency} [\si{\hertz}];
everything else is the derived class's business. The base \cpp{Setup()} schedules
the periodic step and nothing more:

\begin{lstlisting}[language=cpp, caption={The \cpp{Application} lifecycle.}, label={lst:sim-app-lifecycle}]
// cpp/lupnt/applications/application.cc
Application::Application(Config& config) {
  name_ = config["name"] ? config["name"].as<std::string>() : GetId();
  config_ = config;
  if (config["frequency"]) frequency_ = config["frequency"].as<Real>();
}

void Application::Setup() {
  if (frequency_ > 0.0) {
    agent_->GetSimulation()->Schedule(
        0.0, [this](Real t) { Step(t); }, frequency_, Event::Priority::APPLICATION);
  } else {
    Logger::Warn(fmt::format("No frequency set for {}", name_), "Agent");
  }
}

virtual void Step(Real t) = 0;     // pure virtual: one predict/update cycle
virtual void Log(Real t);          // called from Agent::Log
\end{lstlisting}

So the three overridable virtuals are \cpp{Setup()} (resolve target agents, reach
the environment, seed the filter state and covariance, then call the base to keep
the scheduling), the pure-virtual \cpp{Step(Real t)} (one predict/update cycle),
and \cpp{Log(Real t)} (record results; \cpp{Agent::Log} forwards to it).

Logic can also be composed \emph{within} one application through the sub-application
pattern: \cpp{LunaNetSatApp} holds a \cpp{std::vector<Ptr<LunaNetSubApp>>} and fans
its \cpp{Step} out to each sub-application (\cpp{IslOdtsApp},
\cpp{EphemerisGenApp}).

\begin{table}[H]
  \centering
  \caption{Application classes registered with \cpp{ApplicationFactory}.}
  \label{tab:sim-app-classes}
  \begin{tabularx}{\textwidth}{@{}l L@{}}
    \toprule
    \code{class:} & Role \\
    \midrule
    \cpp{GroundStationTrackingApp} & Sensor on each \cpp{GroundStation}:
      visibility-gated range and range-rate, with optional troposphere,
      ionosphere, Shapiro, and solid Earth tide corrections, pushed to a manager. \\
    \cpp{GroundStationManagerApp} & Centralized batch (weighted least squares) plus
      square-root information filter (SRIF) and smoother, with configurable filter
      dynamics. \\
    \cpp{LunarGnssOdtsApp} & GNSS ODTS; calls the Monte-Carlo engine from its
      \cpp{Step}. \\
    \cpp{SatelliteOdtsApp} & Per-satellite onboard Schmidt EKF for the distributed
      ISL scenario; exchanges consider-states over publish--subscribe. \\
    \cpp{StationBeaconSensor} & Surface-station sensor pushing
      \cpp{IslStationObs} records to the ground filter. \\
    \cpp{GroundOdtsApp} & Centralized station-only EKF over every satellite's joint
      orbit $+$ clock 8-state. \\
    \cpp{EphemerisApp}, \cpp{LunaNetSatApp} & Ephemeris fitting; the latter hosts
      sub-applications (\cpp{IslOdtsApp}, \cpp{EphemerisGenApp}). \\
    \cpp{SurfaceStationApp} & Surface-station logic. \\
    \cpp{SurfaceRoverNavApp} & Rover inertial navigation aided by a
      \cpp{LunarNavConstellation}. \\
    \cpp{LanderGncApp} & Guidance / control: owns the powered-descent truth
      trajectory and writes the lander state. \\
    \cpp{LanderNavApp} & Error-state inertial navigation MEKF fusing IMU,
      altimeter, crater bearings, and LunaNet pseudoranges. \\
    \bottomrule
  \end{tabularx}
\end{table}

\subsection{Device, Measurement, Dynamics, State}
\label{sec:sim-models}

\paragraph{Device.} Created from the agent's \code{devices:} block, stored in
\cpp{Agent::devices\_} under \code{<agent>/<device>}, and stepped independently on
its own \code{DEVICE}-priority events, so its internal state is current when an
application reads it with \cpp{GetDevice(name)} (which accepts the bare or the
prefixed name). Registered devices: \cpp{Clock}, \cpp{Imu}, \cpp{Camera},
\cpp{Transmitter}, \cpp{Receiver}, \cpp{Transponder}, \cpp{GnssTransmitter},
\cpp{GnssReceiver}.

\paragraph{Measurement.} The observable contract is one method:

\begin{equation}
  \label{eq:sim-meas}
  \cpp{MeasData Compute(const State\& x, MatXd* H) const},
\end{equation}
which returns
\begin{equation}
  \label{eq:sim-meas-terms}
  \vec{z} = \vec{h}(\vec{x}), \qquad
  \mat{H} = \frac{\partial \vec{h}}{\partial \vec{x}}, \qquad
  \mat{R} = \operatorname{Cov}\!\left[\vec{z}\right].
\end{equation}

\cpp{MeasData} carries \code{timestamp}, the stacked observable \code{value}
($\vec{z}$), and \code{covariance} ($\mat{R}$, sized $n_z \times n_z$, diagonal by
convention). \code{H} is filled only when the caller passes a non-null pointer.
Derive from \cpp{MeasurementClone<Derived>}, a CRTP helper that supplies the
required polymorphic \cpp{Clone()}. \cpp{CreateFunction()} adapts the model to a
\cpp{FilterMeasurementFunction} suitable for
\cpp{Filter::SetMeasurementFunction} by capturing a \cpp{Clone()} of itself:

\begin{lstlisting}[language=cpp, caption={Adapting a measurement model to a filter.}, label={lst:sim-meas-fn}]
// cpp/lupnt/measurements/measurement.h :: Measurement::CreateFunction
virtual FilterMeasurementFunction CreateFunction() const {
  Ptr<Measurement> self = Clone();
  return [self](const State& x, MatXd* H, MatXd* R) -> VecXd {
    MeasData md = self->Compute(x, H);
    if (R != nullptr) *R = md.covariance;
    return md.value;
  };
}
\end{lstlisting}

Error-state (inertial navigation) observables use the parallel
\cpp{ErrorStateMeasurement} base instead. Its Jacobian is taken with respect to
the navigation \emph{error} state
$\delta\vec{x} = [\delta\vec{r}, \delta\vec{v}, \delta\vec{\theta},
\delta\vec{b}_a, \delta\vec{b}_g, \delta b, \delta d]$ and depends on the nominal
attitude, so the model consumes a \cpp{NavErrorContext}
($\vec{r}$, $\vec{v}$, $\mat{R}_{b \to n}$, clock bias and drift, error-state
size) rather than a flat \cpp{State}, and there is no \cpp{Filter} integration ---
the hosting application drives the Joseph-form update directly.

\paragraph{Dynamics.} The \cpp{Propagate} family is the contract. The
state-transition-matrix overload is what filters use for covariance propagation;
if a subclass does not provide one, the base computes it by automatic
differentiation. Concrete classes registered with \cpp{DynamicsFactory}:
\cpp{NBodyDynamics}, \cpp{JointOrbitClockDynamics}, \cpp{StaticDynamics},
\cpp{SurfaceDynamics2D}; \cpp{ClockDynamics} has its own registry, and
\cpp{FixedAttitudeDynamics} / \cpp{FixedPointingDynamics} register as
\cpp{AttitudeDynamics}.

\paragraph{State.} A labeled vector (\cpp{Cart6}, \cpp{ClassicalOE},
\cpp{ClockState2}, \cpp{ImuState}, \dots). \cpp{JointState} concatenates
(\cpp{State}, \cpp{Dynamics}, \cpp{ParamState}, estimation types, $\tau_q$) triples
and exposes the combined dynamics, state-transition-matrix, and process-noise
functions the \cpp{Filter} consumes --- this is how orbit $+$ clock ($+$ estimated
or consider parameters) become one filter state.

\section{Per-step data flow}
\label{sec:sim-dataflow}

At a given epoch the engine runs all due callbacks, and the following work
happens (\cref{fig:sim-object-model}):

\begin{enumerate}[leftmargin=*]
  \item \cpp{AgentWithDynamics::Step} advances each physical agent's \emph{truth}
        state by \cpp{Propagate}, then logs it.
  \item \cpp{Device::Step} advances clocks, IMUs, cameras, and radios.
  \item \cpp{Application::Step} reads other agents' truth via \cpp{GetStateAt(t)}
        or \cpp{World::GetStateAt(name, t)} and its own devices via
        \cpp{GetDevice(...)}, forms one or more \cpp{Measurement}s
        ($\vec{z}, \mat{H}, \mat{R}$), and runs a filter predict (using the
        \cpp{Dynamics} / \cpp{JointState} state-transition matrix and process
        noise) followed by an update.
  \item Results are logged through \cpp{Application::Log}, reached from
        \cpp{Agent::Log}. Periodic events reschedule themselves.
\end{enumerate}

In pseudocode, an agent step and a navigation application step are:

\begin{lstlisting}[language=cpp, caption={Truth propagation and one navigation cycle (pseudocode).}, label={lst:sim-step-pseudocode}]
// AgentWithDynamics::Step(t)          (priority AGENT) -- advance the truth state
state_, attitude_ = dynamics_.Propagate(state_, time_, t, control_);
time_ = t;  Log(t);

// A navigation Application::Step(t)   (priority APPLICATION)
for (tx : visible_transmitters(t)) {          // geometry via tx.GetStateAt(t)
  z, H, R = measurement(tx, receiver = agent_).Compute(x_hat, &H);
  stack(z, H, R);
}
filter.Predict(t);                            // x_hat, P via Dynamics/JointState STM + Q
filter.Update(z, H, R);                       // innovation, gain, covariance update
Log(t);
\end{lstlisting}

\section{Configuration and the asset factory}
\label{sec:sim-factory}

Classes are instantiated by name through \cpp{AssetFactory<Base, Config\&>}
(\file{cpp/lupnt/core/asset_factory.h}). A class opts in with one macro at file
scope:

\begin{lstlisting}[language=cpp, caption={Factory registration.}, label={lst:sim-register}]
// cpp/lupnt/core/asset_factory.h :: REGISTER_FACTORY_CLASS(BASE, DERIVED)
// registers a creator  [](Config& config) { return std::make_shared<DERIVED>(config); }
// under the string "DERIVED", via a file-scope static registrar object.

REGISTER_FACTORY_CLASS(Application, MyApp)    // at file scope in my_app.cc
\end{lstlisting}

The YAML \code{class:} value must match that string exactly; if it does not, the
factory aborts with a message listing every registered class name. Typed aliases
exist for each base: \cpp{AgentFactory}, \cpp{DeviceFactory},
\cpp{ApplicationFactory}, \cpp{DynamicsFactory}, \cpp{ChannelFactory}. Each
registry is defined in exactly one translation unit and explicitly instantiated,
so a single registry is shared across library boundaries.

\subsection{The scenario schema}
\label{sec:sim-schema}

\cpp{Simulation::Setup} reads the scenario-level keys in
\cref{tab:sim-top-keys}, in that order, then hands each \code{agents:} entry to
the factory named by its \code{class:}. The agent constructor in turn builds its
\code{dynamics:}, \code{attitude\_dynamics:}, \code{initial\_state:},
\code{devices:}, and \code{application:} (or \code{applications:}) sub-blocks
through the matching factories. Names are made unique by prefixing:
a channel, agent, or constellation named \code{X} in a simulation named \code{S}
is registered under \code{S/X}; a device \code{d} on agent \code{A} under
\code{A/d}; an application under \code{A/<name>}, defaulting to
\code{A/application} for the singular form and to \code{A/<class>} for each entry
of an \code{applications:} list.

\begin{table}[H]
  \centering
  \caption{Top-level scenario keys parsed by \cpp{Simulation::Setup}.}
  \label{tab:sim-top-keys}
  \begin{tabularx}{\textwidth}{@{}l L@{}}
    \toprule
    Key & Meaning \\
    \midrule
    \code{name} & Scenario name; also the output directory and the
      \code{<name>.h5} log file, and the registry prefix. \\
    \code{log\_level} & One of the \cpp{Logger::LogLevel} enum names
      (e.g.\ \code{DEBUG}, \code{INFO}, \code{WARNING}). \\
    \code{cesium\_port} & If present, starts a \cpp{CesiumViewer} on that port. \\
    \code{epoch} & Gregorian string, e.g.\ \code{"2027-03-01T00:00:00"}; sets the
      global \lupnt{} epoch for $t = 0$. \\
    \code{duration} & \code{HH:MM:SS(.sss)}; parsed to seconds and used as the
      event-loop horizon. \\
    \code{world} & The shared environment (\cref{tab:sim-world-keys}). \\
    \code{channels} & Map of \cpp{Channel} configs. \\
    \code{agents} & Map of agent configs, each with a \code{class:}. \\
    \code{constellations} & Map of \cpp{Constellation} configs. \\
    \code{seed} & Monte-Carlo seed, propagated as the default to every agent and
      its application. \\
    \bottomrule
  \end{tabularx}
\end{table}

An abridged scenario, from \file{configs/ground_station_odts.yaml}:

\begin{lstlisting}[language=yamlcfg, caption={A complete scenario skeleton, abridged from \texttt{configs/ground\_station\_odts.yaml}.}, label={lst:sim-yaml-gsodts}]
# configs/ground_station_odts.yaml
name: ground_station_odts
epoch: "2026-01-01T00:00:00"
duration: "36:00:00"          # 36 h (~2.7 ELFO orbits)
log_level: INFO

world:                        # shared, read-only environment
  frame: MOON_CI
  force_model:                # 2x2 Moon + Earth/Sun + SRP
    integrator: RKF45
    dt: 10.0
    max_iter: 1000
    abstol: 1.0e-6
    reltol: 1.0e-10
    autodiff: true            # analytic STM for the estimator
    mass: 10.0                # [kg]   } SRP: CR * area / mass
    area: 0.02                # [m^2]  }
    CR: 1.8
    bodies:
      - MOON: { n: 2, m: 2 }
      - EARTH: {}
      - SUN: {}

agents:
  sat:                              # truth target
    class: Satellite
    frequency: 0.0033333333333333   # 1/300 Hz truth logging
    # No `dynamics:` block -> inherits world.force_model (autodiff forced off).
    initial_state:
      class: ClassicalOE
      frame: MOON_OP
      a: 6541400.0                  # [m]
      e: 0.6
      i: 56.2                       # [deg]
      Omega: 0.0
      omega: 90.0
      M: 0.0

  gs_manager:                       # centralized OD estimator
    class: GroundStationManager
    application:
      class: GroundStationManagerApp
      target: sat
      obs_interval_s: 300.0
      seed: 42
      initial_position_sigma_m: 2000.0
      initial_velocity_sigma_mps: 0.5
      batch_max_iterations: 10
      batch_convergence_tol: 1.0e-3
      run_srif: true
      srif_use_process_noise: true
      srif_accel_psd: 3.0e-13

  DSS14:                            # a tracking station (sensor)
    class: GroundStation
    latitude_deg: 35.426456
    longitude_deg: 243.110461
    altitude_m: 1001.39
    application:
      class: GroundStationTrackingApp
      frequency: 0.0033333333333333
      target: sat                   # resolved via GetSimulation()->GetAgent("sat")
      manager: gs_manager           # where observations are pushed
      elevation_mask_deg: 10.0
      use_range: true
      use_range_rate: true
      range_sigma_m: 10.0
      range_rate_sigma_mps: 1.0e-3
      apply_troposphere: true       # Saastamoinen zenith delay + 1/sin(el) mapping
      apply_ionosphere: true        # thin-shell VTEC delay 40.3*STEC/f^2
      apply_shapiro: true           # relativistic light-time delay
      apply_solid_earth_tide: true  # degree-2 station displacement
\end{lstlisting}

\begin{lupntwarning}
\code{yaml-cpp} does not support YAML merge keys (\code{<<:}), so repeated station
blocks must be written out in full. Anchors and aliases (\code{\&id}, \code{*id})
\emph{are} supported and are used in
\file{configs/isl_odts_distributed.yaml} to share the station list across
satellites.
\end{lupntwarning}

\subsection{One force-model schema everywhere}
\label{sec:sim-force-model}

Every place that specifies an orbit force model --- the shared
\code{world: force\_model:}, a physical agent's \code{dynamics:}, and an
estimator's filter force model --- uses the \emph{same} block, parsed by
\cpp{ParseForceModelSpec} (\file{cpp/lupnt/dynamics/numerical_orbit_dynamics.cc}).
There is no separate \code{moon\_gravity\_degree\_*} / \code{include\_earth}
dialect. The keys are:

\begin{table}[H]
  \centering
  \caption{The shared force-model block.}
  \label{tab:sim-force-model}
  \begin{tabularx}{\textwidth}{@{}l L@{}}
    \toprule
    Key & Meaning \\
    \midrule
    \code{bodies} & Sequence of single-key maps \code{BODY: \{ n, m \}}. For
      \code{MOON}, \code{n}/\code{m} are the spherical-harmonic degree and order
      (default $0$, i.e.\ point mass); \code{EARTH} and \code{SUN} are third-body
      point masses and take \code{\{\}}. \\
    \code{relativity} & Enable the relativistic correction. The spelling
      \code{use\_relativity} is accepted as an alias. \\
    \code{CR}, \code{area}, \code{mass} & Solar-radiation-pressure ballistic
      coefficient $C_R A / m$; SRP is enabled only when \code{CR} is present. \\
    \code{CD} & Drag coefficient (with \code{area} and \code{mass}). \\
    \code{integrator}, \code{dt}, \code{abstol}, \code{reltol}, \code{max\_iter} &
      Integrator selection and step/tolerance controls. \\
    \code{autodiff} & Whether the model carries an analytic state-transition
      matrix. Estimators need \code{true}; truth propagation does not. \\
    \code{frame} & Reference frame; inherited from \code{world.frame} when
      omitted in an injected block. \\
    \code{units} & Optional non-dimensional unit system. \\
    \bottomrule
  \end{tabularx}
\end{table}

\subsection{Different dynamics for truth and filter}
\label{sec:sim-truth-filter}

Truth and filter dynamics are independent knobs, so a realistic orbit-determination
study can run the estimator against a deliberately mismodeled truth.

\emph{Truth} is whatever propagates the physical agent: a self-propagating
\cpp{AgentWithDynamics} integrates its own \code{dynamics:} block (or the injected
\code{world.force\_model}), and an application reads that truth through
\cpp{GetStateAt(t)} or \cpp{World::GetStateAt(name, t)}.

\emph{Filter} dynamics are built \emph{inside} the estimator \cpp{Application} ---
either by reusing \cpp{World::MakeDynamics()} or from application-specific config.
Whatever the filter uses also supplies its state-transition matrix, so keep
\code{autodiff: true} on that model.

\emph{Same model (no mismatch).} Omit the target agent's \code{dynamics:} block so
it inherits \code{world: force\_model:}, and have the application build its filter
from \cpp{World::MakeDynamics()}. This is exactly what
\file{configs/ground_station_odts.yaml} and \file{configs/sat_bearing_odts.yaml}
do.

\emph{Different models (mismatch).} Configure the two sides separately. In the
distributed ISL scenario the truth is propagated from \code{world: force\_model:}
--- inherited by every \cpp{Spacecraft} --- at $16 \times 16$ lunar gravity, while
each onboard \cpp{SatelliteOdtsApp} runs its EKF at $8 \times 8$:

\begin{lstlisting}[language=yamlcfg, caption={Deliberate truth/filter force-model mismatch.}, label={lst:sim-yaml-mismatch}]
# configs/isl_odts_distributed.yaml
world:
  force_model:                            # truth propagated at 16x16 lunar gravity
    integrator: RKF45
    dt: 60.0
    autodiff: false
    bodies:
      - MOON: { n: 16, m: 16 }
      - EARTH: {}
      - SUN: {}

# ... every Spacecraft inherits that truth force model; the onboard app filters
# coarser using the SAME force_model schema:
application:
  class: SatelliteOdtsApp
  force_model:                            # EKF runs at 8x8 -> intentional mismatch
    bodies:
      - MOON: { n: 8, m: 8 }
      - EARTH: {}
      - SUN: {}
    relativity: true
  integration_step_s: 60.0
  process_accel_sigma_mps2: 1.0e-08
\end{lstlisting}

More generally: give the truth agent a high-fidelity \code{dynamics:} block (more
third bodies, higher gravity degree and order, SRP, drag) and point the filter at
a coarser model. An estimator application can also expose the filter model as
config: \cpp{GroundStationManagerApp} accepts an optional \code{filter\_dynamics:}
block applied to \emph{both} its batch and its sequential SRIF/smoother filters,
or \code{batch\_dynamics:} / \code{sequential\_dynamics:} to set each
independently. Absent those, both inherit \code{world: force\_model:} --- truth
equals filter, the default. In a custom application, simply build two
\cpp{Dynamics} objects instead of sharing one.

\section{Authoring a new scenario}
\label{sec:sim-authoring}

The reference template is \file{configs/ground_station_odts.yaml} together with
\file{python/examples/ex7_run_odts.py}. The steps, in increasing order of effort:

\begin{enumerate}[leftmargin=*]
  \item \textbf{Add the shared environment.} Write a \code{world:} block
    (\code{frame} $+$ \code{force\_model}, or a \code{gravity:} / \code{dem:}
    block for surface scenarios). Both truth and estimator dynamics derive from
    it.
  \item \textbf{Reuse existing classes.} List \code{agents:} with
    already-registered \code{class:} values
    (\cref{tab:sim-agent-classes}), each with an \code{initial\_state:}, an
    optional \code{dynamics:} block, optional \code{devices:}, and an
    \code{application:} block. Launch with \py{pnt.Simulation(config)}. Nothing
    else is needed if the physics and the observables already exist.
  \item \textbf{New Application} --- the most common extension. Subclass
    \cpp{Application}, implement \cpp{MyApp(Config\&)}, \cpp{Setup()}, and
    \cpp{Step(Real)}; add \cpp{REGISTER\_FACTORY\_CLASS(Application, MyApp)}.
    Reference: \file{cpp/lupnt/applications/ground_station/ground_station_manager_app.cc}.
  \item \textbf{New Measurement} --- only for a new observable. Subclass
    \cpp{MeasurementClone<MyMeas>} (or \cpp{ErrorStateMeasurement}) and implement
    \cpp{Compute(const State\&, MatXd*)}. Wire it into a filter with
    \cpp{CreateFunction()}. Reference:
    \file{cpp/lupnt/measurements/crosslink_measurement.cc}.
  \item \textbf{New Dynamics} --- only for new physics. Subclass
    \cpp{NumericalDynamics} (implementing \cpp{ComputeRates}, \cpp{Propagate},
    \cpp{GetStateType}) or \cpp{Dynamics} directly, and add
    \cpp{REGISTER\_FACTORY\_CLASS(Dynamics, MyDyn)}. The state-transition matrix is
    supplied by automatic differentiation if you do not provide one.
  \item \textbf{New Agent} --- only for a new platform kind. Subclass
    \cpp{AgentWithDynamics}, implement \cpp{GetStateAt}, optionally override
    \cpp{Propagate} / \cpp{Setup}, and add
    \cpp{REGISTER\_FACTORY\_CLASS(Agent, MyAgent)}. Reference:
    \file{cpp/lupnt/agents/satellite.cc}, \file{cpp/lupnt/agents/rover.cc}.
\end{enumerate}

The skeleton of the two most common C++ extensions:

\begin{lstlisting}[language=cpp, caption={A new application and a new observable in C++.}, label={lst:sim-cpp-skeleton}]
// --- A new navigation application ---------------------------------------
class MyOdtsApp : public Application {
 public:
  MyOdtsApp(Config& cfg) : Application(cfg) {
    target_ = cfg["target"].as<std::string>();       // agent name to track
    sigma_  = cfg["range_sigma_m"].as<double>();
  }
  void Setup() override {
    Application::Setup();                            // schedules Step at frequency_
    target_agent_ = GetAgent()->GetSimulation()->GetAgent(target_);
    dyn_ = GetAgent()->GetWorld()->MakeDynamics();    // filter model, STM enabled
    // seed filter_ state x_hat_ and covariance P_ here
  }
  void Step(Real t) override {
    Cart6 x_tx = target_agent_->GetStateAt(t);       // truth geometry
    MatXd H; MyMeasurement meas(BuildConfig(x_tx));
    MeasData zHR = meas.Compute(x_hat_, &H);         // z = h(x), R, and H
    filter_.Predict(t);                              // STM + process noise
    filter_.Update(zHR.value, H, zHR.covariance);    // innovation, gain, update
    Log(t);
  }
};
REGISTER_FACTORY_CLASS(Application, MyOdtsApp)       // YAML: class: MyOdtsApp

// --- A new observable ---------------------------------------------------
class MyMeasurement : public MeasurementClone<MyMeasurement> {
 public:
  MeasData Compute(const State& x, MatXd* H = nullptr) const override {
    VecXd z(1);   z(0) = /* h(x): range, Doppler, ... */;
    if (H) { H->resize(1, x.size()); *H = /* dh/dx */; }
    return MeasData{timestamp_, z, R_};              // value, covariance
  }
};
\end{lstlisting}

\begin{lupntwarning}
\cpp{Simulation(config)} calls \cpp{Setup()} inside its constructor, so every
agent and application is fully built and set up before \cpp{Run()} is reachable.
Any state that a host must inject \emph{after} construction (for example the
lander guidance application's reference trajectory) should therefore be applied
through a lazy \cpp{Initialize} inside the application's first \cpp{Step}.
\end{lupntwarning}

\section{Working in Python}
\label{sec:sim-python}

The whole simulation lifecycle is driven from Python. A scenario is just a
\py{dict} (or a loaded YAML document), so it can be built, mutated, and swept
programmatically:

\begin{lstlisting}[language=py, caption={Composing, mutating, and running a scenario from Python.}, label={lst:sim-py-run}]
# python/examples/ex7_run_odts.py (pattern)
import copy, yaml, pylupnt as pnt

cfg = yaml.safe_load(open("configs/ground_station_odts.yaml"))
cfg["agents"]["gs_manager"]["application"]["srif_accel_psd"] = 1e-13

sim = pnt.Simulation(cfg)     # builds World + agents + apps from the `class:` keys
sim.run()                     # runs the event loop
app = sim.get_agent("gs_manager").get_application()   # downcasts to the concrete app
\end{lstlisting}

\cpp{pnt.Simulation} accepts either a path string or a config dict. The bound
surface is deliberately small: \py{run()}, \py{get\_agent(name)},
\py{get\_world()}, \py{get\_duration()}, \py{get\_time()} on the simulation;
\py{get\_name()}, \py{get\_state\_at(t)}, \py{get\_application()},
\py{get\_applications()}, \py{get\_application\_by\_name(name)}, \py{get\_world()} on
an agent; and \py{get\_epoch()}, \py{get\_frame()}, \py{has\_force\_model()},
\py{get\_gm()}, \py{gravity(r)}, the DEM accessors, and \py{get\_state\_at(name, t)}
on the world. Each concrete application binds its own result accessors in its own
binding file (\file{python/bindings/py_ground_station_odts.cc},
\file{python/bindings/py_isl_odts.cc}, and so on).

\subsection{Authoring an Application in Python}
\label{sec:sim-py-application}

\cpp{Application}, \cpp{Agent}, and \cpp{Measurement} are exposed as polymorphic
bases with \code{pybind11} trampolines (\cpp{PyApplication}, \cpp{PyAgent},
\cpp{PyMeasurement} in \file{python/bindings/py_simulation.cc}), so a pure-Python
subclass can be driven by the C++ engine with no C++ build.
\Cref{tab:sim-trampolines} gives the dispatch contract.

\begin{table}[H]
  \centering
  \caption{Python override points exposed by the trampolines in
           \texttt{python/bindings/py\_simulation.cc}.}
  \label{tab:sim-trampolines}
  \begin{tabularx}{\textwidth}{@{}l l L@{}}
    \toprule
    Base & Python method & Contract \\
    \midrule
    \py{pnt.Application} & \py{step(self, t)} & \textbf{Required.} $t$ arrives as
      a plain Python \py{float} (not an autodiff \cpp{Real}), so ordinary
      arithmetic works. \\
    & \py{setup(self)} & Optional. Call
      \py{pnt.Application.setup(self)} to keep the base's frequency-based
      scheduling of \py{step}. \\
    & \py{log(self, t)} & Optional; falls back to the C++ base. \\
    \py{pnt.Agent} & \py{get\_state\_at(self, t)} & \textbf{Required.} Must return a
      numpy 6-vector $[\vec{r}; \vec{v}]$; the trampoline checks the size. \\
    & \py{setup}, \py{step}, \py{log} & Optional; fall back to the C++ base. \\
    \py{pnt.Measurement} & \py{compute(self, x)} & \textbf{Required.} Returns the
      tuple $(\vec{z}, \mat{H} = \partial \vec{h}/\partial \vec{x}, \mat{R})$, all
      numpy. \py{evaluate(x)} runs it through the C++ base --- the same path a
      filter uses. \\
    \bottomrule
  \end{tabularx}
\end{table}

Register the class with \py{pnt.register\_application(name, cls)} and reference it
by \code{class: name} in a config. The factory constructs it as
\py{cls(config\_dict)}, where the dict is the \code{application:} block converted
recursively from YAML (integers, floats --- including scientific notation such as
\code{1.0e-8} --- booleans, strings, lists, dicts). The returned Python object is
kept alive for the simulation's lifetime and
\py{sim.get\_agent(...).get\_application()} hands back the \emph{same} object, so
results stored on \py{self} survive the run.

\begin{lstlisting}[language=py, caption={An estimator application authored entirely in Python.}, label={lst:sim-py-app}]
# python/examples/ex17_python_new_sim_example.ipynb
import numpy as np, yaml, pylupnt as pnt

def make_dyn():
    d = pnt.NBodyDynamics()
    for b in (pnt.Body.Moon(8, 8), pnt.Body.Earth(), pnt.Body.Sun()):
        d.add_body(b)
    d.set_frame(pnt.Frame.MOON_CI)
    d.set_integrator(pnt.IntegratorType.RKF45)
    d.set_autodiff(True)          # analytic STM for the filter
    return d

class PyAnglesOdtsApp(pnt.Application):
    def __init__(self, config):                      # config = the `application:` block
        pnt.Application.__init__(self)
        self.target = config["target"]
        self.set_frequency(config.get("frequency", 1.0 / 60))
        self.sigma = np.radians(config.get("angle_sigma_arcsec", 2.0) / 3600.0)
        self.p0 = config.get("initial_position_sigma_m", 300.0)
        self.v0 = config.get("initial_velocity_sigma_mps", 0.3)
        self.qa = config.get("process_accel_sigma_mps2", 1e-6)
        self.rng = np.random.default_rng(config.get("seed", 42))
        self.times, self.est, self.truth, self.sig = [], [], [], []
        self._init = False

    def setup(self):
        pnt.Application.setup(self)   # base schedules step() at get_frequency()
        self.dyn = make_dyn()
        self.meas = BearingMeasurement(self.sigma)

    def step(self, t):
        if not self._init:                                   # seed once
            self.obs = np.asarray(self.get_agent().get_state_at(0.0))
            self.tgt = np.asarray(
                self.get_agent().get_world().get_state_at(self.target, 0.0))
            self.x = self.obs + np.r_[self.rng.normal(0, self.p0, 3),
                                      self.rng.normal(0, self.v0, 3)]
            self.P = np.diag(np.r_[[self.p0**2] * 3, [self.v0**2] * 3])
            self._init, self._t = True, t
            return
        dt = t - self._t
        xf, F = self.dyn.propagate_stm(self.x, float(self._t), float(t))
        self.x = np.asarray(xf).ravel(); F = np.asarray(F); self._t = t
        Q = self.qa**2 * np.block([[dt**3 / 3 * np.eye(3), dt**2 / 2 * np.eye(3)],
                                   [dt**2 / 2 * np.eye(3), dt * np.eye(3)]])
        self.P = F @ self.P @ F.T + Q                        # predict
        self.meas.target = self.tgt[:3]
        z = self.meas.evaluate(self.obs)[0] + self.rng.normal(0, self.sigma, 3)
        u, H, R = self.meas.evaluate(self.x)                 # predict the observable
        S = H @ self.P @ H.T + R
        K = self.P @ H.T @ np.linalg.inv(S)
        self.x = self.x + K @ (z - u)                        # update
        IKH = np.eye(6) - K @ H
        self.P = IKH @ self.P @ IKH.T + K @ R @ K.T          # Joseph form
        self.times.append(t); self.est.append(self.x.copy())

pnt.register_application("PyAnglesOdtsApp", PyAnglesOdtsApp)
cfg = yaml.safe_load(open("configs/sat_bearing_odts.yaml"))
sim = pnt.Simulation(cfg); sim.run()
est = np.array(sim.get_agent("observer").get_application().est)   # results off `self`
\end{lstlisting}

The scenario it runs is \file{configs/sat_bearing_odts.yaml}, an ordinary config
whose only unusual feature is that its \code{class:} names a Python type:

\begin{lstlisting}[language=yamlcfg, caption={A config whose application class is authored in Python.}, label={lst:sim-yaml-bearing}]
# configs/sat_bearing_odts.yaml (abridged)
observer:
  class: Spacecraft
  frequency: 0.016666666666666666
  initial_state:
    class: ClassicalOE
    frame: MOON_OP
    a: 1837400.0                    # [m] ~100 km altitude circular
    e: 0.001
    i: 30.0
    Omega: 0.0
    omega: 0.0
    M: 0.0
  application:
    class: PyAnglesOdtsApp          # registered by pnt.register_application
    target: target
    frequency: 0.016666666666666666
    dt_s: 60.0
    angle_sigma_arcsec: 2.0         # star-tracker class bearing noise
    seed: 42
    force_model:                    # filter force model, same schema as world
      bodies:
        - MOON: { n: 8, m: 8 }
        - EARTH: {}
        - SUN: {}
      relativity: false
    process_accel_sigma_mps2: 1.0e-6
    initial_position_sigma_m: 300.0
    initial_velocity_sigma_mps: 0.3
\end{lstlisting}

\subsection{Authoring a Measurement in Python}
\label{sec:sim-py-measurement}

A \py{pnt.Measurement} subclass implements
\py{compute(self, x) -> (z, H, R)}. Because \py{evaluate(x)} routes through the
C++ base --- the same path a \cpp{Filter} uses --- the \emph{same} Python class
both \emph{generates} an observation (apply it to a truth state and add noise) and
\emph{predicts} one (apply it to the filter state). The bearing model of
Example~17 is the unit line of sight $\uvec{u}$ from the observer at $\vec{r}$ to a
known target at $\vec{p}$:

\begin{align}
  \vec{d} &= \vec{p} - \vec{r}, \qquad \rho = \lVert \vec{d} \rVert, \qquad
  \uvec{u} = \vec{d}/\rho, \label{eq:sim-bearing}\\
  \frac{\partial \uvec{u}}{\partial \vec{r}}
    &= -\frac{1}{\rho}\left(\mat{I}_3 - \uvec{u}\,\uvec{u}\transpose\right),
  \qquad
  \frac{\partial \uvec{u}}{\partial \vec{v}} = \mat{0}_{3\times 3},
  \qquad
  \mat{R} = \sigma^2 \mat{I}_3. \label{eq:sim-bearing-jac}
\end{align}

\begin{lstlisting}[language=py, caption={A measurement model authored in Python.}, label={lst:sim-py-meas}]
# python/examples/ex17_python_new_sim_example.ipynb
class BearingMeasurement(pnt.Measurement):
    "Unit line-of-sight (bearing) to a known target position."

    def __init__(self, sigma_rad):
        pnt.Measurement.__init__(self)
        self.sigma = sigma_rad
        self.target = np.zeros(3)    # set per-epoch to the target truth position

    def compute(self, x):            # -> (z = h(x), H = dh/dx, R)
        d = self.target - x[:3]
        rn = np.linalg.norm(d)
        u = d / rn
        H = np.zeros((3, 6))
        H[:, :3] = -(np.eye(3) - np.outer(u, u)) / rn
        return u, H, self.sigma**2 * np.eye(3)

m = BearingMeasurement(np.radians(2 / 3600))
m.target = tgt_truth[:3]
z = m.evaluate(obs_truth)[0] + rng.normal(0, m.sigma, 3)   # GENERATE from truth
u, Hx, R = m.evaluate(x_hat)                               # PREDICT at the filter state
\end{lstlisting}

\subsection{Authoring an Agent in Python}
\label{sec:sim-py-agent}

A platform whose truth trajectory is defined in Python --- an analytic ephemeris, a
scripted path --- subclasses \py{pnt.Agent}, implements \py{get\_state\_at(self, t)},
and is registered with \py{pnt.register\_agent(name, cls)}:

\begin{lstlisting}[language=py, caption={A physical platform whose truth is an analytic Python ephemeris.}, label={lst:sim-py-agent}]
class PyEphemeris(pnt.Agent):
    def __init__(self, config):
        pnt.Agent.__init__(self)
        self.a = config["radius_m"]
        self.w = (pnt.GM_MOON / self.a**3) ** 0.5
    def get_state_at(self, t):                 # must return a numpy [r; v] 6-vector
        c, s = np.cos(self.w * t), np.sin(self.w * t)
        return np.array([self.a * c, self.a * s, 0.0,
                         -self.a * self.w * s, self.a * self.w * c, 0.0])

pnt.register_agent("PyEphemeris", PyEphemeris)
# config:  {class: PyEphemeris, radius_m: 2.0e6}
\end{lstlisting}

A complete, converging example --- the angles-only orbit-determination EKF of
\cref{lst:sim-py-app} and \cref{lst:sim-py-meas}, reaching the metre level over a
\SI{6}{\hour} arc from a \SI{300}{\meter} a-priori with \SI{2}{arcsecond} bearings
--- is Example~17, \file{python/examples/ex17_python_new_sim_example.ipynb}.

\subsection{Monte-Carlo trials}
\label{sec:sim-montecarlo}

Monte-Carlo settings live at the \emph{simulation} level, not on an agent or an
application: a trial is one independent \py{pnt.Simulation} run built from a deep
copy of the config with a distinct top-level \code{seed:}, which
\cpp{Simulation::Setup} propagates as the default to every agent and every
application. \py{pnt.run\_monte\_carlo} (\file{python/pylupnt/montecarlo.py})
parallelizes across \emph{processes} --- the C++ engine stays serial, and a
\code{spawn} context gives each worker a clean interpreter so the thread-local C++
random engine starts fresh:

\begin{lstlisting}[language=py, caption={Parallel Monte-Carlo over a scenario config.}, label={lst:sim-py-montecarlo}]
# python/pylupnt/montecarlo.py
import pylupnt as pnt

def final_error(sim):                     # must be a top-level (picklable) function
    app = sim.get_agent("gs_manager").get_application()
    return float(app.est_central()[-1, 0])

summaries = pnt.run_monte_carlo(config, n_runs=16, extract=final_error,
                                n_workers=8, seed0=0, output_root="output/mc")
# trial i uses seed = seed0 + i; every nested `output_dir` key is redirected to
# <output_root>/run_<i> so parallel trials never overwrite each other.
\end{lstlisting}

For a scenario with an expensive \emph{shared} precompute, the GNSS ODTS engine
also offers an in-configuration \code{monte\_carlo\_runs:} knob (one precompute,
then \code{seed = base\_seed + i} per run); see
\file{python/examples/ex6_gnss_odts.ipynb}.

\section{Model boundaries}
\label{sec:sim-boundaries}

This section states what the YAML front end can express on its own, what still
requires code, and which behaviours are approximations of the engine rather than
of the physics.

\subsection{What the YAML front end supports}

\begin{itemize}[leftmargin=*]
  \item The full scenario skeleton: \code{name}, \code{epoch}, \code{duration},
    \code{log\_level}, \code{cesium\_port}, \code{seed}, \code{world},
    \code{channels}, \code{agents}, \code{constellations}
    (\cref{tab:sim-top-keys}).
  \item Any composition of already-registered agents, devices, dynamics,
    channels, and applications, selected by string (\cref{tab:sim-agent-classes},
    \cref{tab:sim-app-classes}); every constructor knob of those classes is a
    config key.
  \item The complete force-model specification (\cref{tab:sim-force-model}),
    used identically for the shared environment, per-agent truth, and per-filter
    fidelity --- so a truth/filter mismatch study is a pure configuration change.
  \item Multi-application platforms (\code{applications:}), constellation
    expansion (\code{walker:} or an explicit \code{satellites:} list), and
    seed-driven Monte-Carlo sweeps.
  \item Any class authored in \emph{Python} and registered with
    \py{pnt.register\_application} or \py{pnt.register\_agent} before
    \py{pnt.Simulation} is constructed --- these are indistinguishable from C++
    classes to the config.
\end{itemize}

\subsection{What still requires code}

\begin{itemize}[leftmargin=*]
  \item \textbf{New physics.} A new \cpp{Dynamics} class must be C++: there is no
    \cpp{PyDynamics} trampoline, so a Python author reuses
    \py{pnt.NBodyDynamics} (built and propagated with state-transition matrix from
    Python) rather than writing a new propagator.
  \item \textbf{New devices and channels.} \cpp{Device} and \cpp{Channel} have no
    Python trampoline either; new hardware models are C++ subclasses with
    \cpp{REGISTER\_FACTORY\_CLASS}.
  \item \textbf{Error-state observables.} \cpp{ErrorStateMeasurement} is not bound
    to Python; only the state-vector \cpp{Measurement} base is. Inertial-navigation
    observables are therefore C++.
  \item \textbf{New filters.} The \cpp{Filter} family (\cpp{EKF}, \cpp{SRIF},
    \cpp{UDUEKF}, \cpp{UDUStochasticCloningEKF}) is registered but not
    Python-subclassable; a Python application implements its own estimator in
    numpy, as \cref{lst:sim-py-app} does.
  \item \textbf{Cross-application data structures.} An application that must hand
    typed records to another (for example \cpp{IslStationObs} from
    \cpp{StationBeaconSensor} to \cpp{GroundOdtsApp}) needs a C++ struct and an
    accessor; the generic escape hatch is the \cpp{Publish} / \cpp{Subscribe} bus,
    which carries a \cpp{std::any}.
  \item \textbf{New initial-state parameterizations.} \cpp{Satellite} accepts only
    a \code{ClassicalOE} block; \cpp{Spacecraft} accepts \code{ClassicalOE} or
    explicit \code{r0\_m}/\code{v0\_mps}. Anything else is a constructor change.
\end{itemize}

\subsection{Engine approximations and expected behaviour}

\begin{itemize}[leftmargin=*]
  \item \textbf{Tie-breaking.} At equal event time the numerically lower priority
    value runs first, so applications may execute before the agents' own
    \cpp{Step} at that epoch. This is safe because \cpp{GetStateAt(t)}
    self-propagates a copy; it is \emph{not} safe for logic that assumes a device
    has already been advanced.
  \item \textbf{Horizon.} \cpp{Run} stops when the queue is empty or the next event
    exceeds \code{duration\_} $+$ \code{EPS}, and a periodic event is not
    rescheduled past that horizon. The final epoch of a run is therefore the last
    multiple of $1/f$ within the duration, not the duration itself.
  \item \textbf{Truth propagation carries no state-transition matrix.} Inherited
    world dynamics are injected with \code{autodiff: false}, and both
    \cpp{Spacecraft} and \cpp{World::MakeTruthDynamics} force it off. An
    application that needs an STM must build its own dynamics with
    \code{autodiff: true}.
  \item \textbf{The \cpp{World} is read-only and never propagates.} It provides
    the environment and the \cpp{GetStateAt} facade; every
    \cpp{AgentWithDynamics} self-propagates. There is no central integrator and no
    global synchronization of agent epochs.
  \item \textbf{Setup runs in the constructor.} There is no window between
    construction and \cpp{Run()} in which the engine will call back into a
    partially built application, so post-construction injection must be deferred
    to the first \cpp{Step}.
  \item \textbf{Determinism.} A run is reproducible for a fixed top-level
    \code{seed:} on a fixed build; parallel Monte-Carlo uses one process per trial
    precisely because the C++ random engine is thread-local, and per-trial
    \code{output\_dir} redirection is required to keep file outputs from
    colliding.
  \item \textbf{Python overhead.} A Python-authored application acquires the GIL on
    every \py{step}, \py{setup}, \py{log}, \py{get\_state\_at}, and \py{compute}
    call. This is acceptable at the \SI{1}{\per\minute} cadences used by the
    orbit-determination examples and is not appropriate for high-rate
    (\si{\kilo\hertz}) inertial mechanization, which stays in C++.
\end{itemize}
